% ===================================================================
%  TABELO N-RO 9 — La monto. — La banurbo. — La arbaro. — La ĉasado.
%  Traduction 2026 d'apres texto/io, controlee sur texto/fr et
%  texto/es. Consigne : voir texto/eo/10-tabelo-01.tex.
% ===================================================================

%%K t09-apar-1 apar
\VUsaut{4.0mm}
\VUtitre{15.0pt}{60}{TABELO N\textsuperscript{o} 9}
\VUsaut{3.0mm}

%%K t09-tit-1 sub
\VUcentre{12.0pt}{0}{\textit{Unua sceno.}}
\VUsaut{3.0mm}

%%K t09-tit-2 sub
\VUpk{11.4pt}{40}{La monto. --- La banurbo.}
\VUsaut{3.0mm}

%%K t09-c1-01-1 p
1. --- De ok tagoj mi estas en Prats. Mi ne povas esprimi la tutan admiron kiun
mi sentis kiam mi estis antaŭ la mirinda \VUgras{montaro} \textsuperscript{(1)}
de la Pireneoj kies \VUgras{kresto} \textsuperscript{(2)} aspektas sur la blua
ĉielo kiel giganta festono.

%%K t09-c1-02-1 p
2. --- Mi ne imagis, ke la pireneaj \VUgras{supraĵoj} \textsuperscript{(3)}
atingas \VUgras{altecon} tiel grandan, kaj mi restis mirigita antaŭ la belegaj
\VUgras{glaĉeroj} \textsuperscript{(5)} kaj la \VUgras{montopintoj}
\textsuperscript{(4)} neĝaj sur kiuj ruliĝas tiel terurigaj \VUgras{lavangoj}.

%%K t09-tit-3 sub
\VUsaut{3.0mm}
\VUpk{10.2pt}{40}{Montopejzaĝo.}
\VUsaut{3.0mm}

%%K t09-c1-03-1 p
3. --- Tuj de la sekva tago de mia alveno, mi iris al la malnova \VUgras{vulkano}
\textsuperscript{(6)} estingita kiu estas apud Prats. Sur la bordoj aridaj kaj
nudigitaj de la \VUgras{kratero}, oni vidas ankoraŭ bone la spurojn de la
paso de la \VUgras{lafo} kiu fluis, plurajn jarcentojn antaŭ nun, sur la
\VUgras{flanko de la monto} \textsuperscript{(7)}. Mi sukcesis trovi, sur unu el
la \VUgras{deklivoj}, kelkajn fragmentojn de tiu lafo.

%%K t09-c1-04-1 p
4. --- Prats estas situita sur la landlimo, kaj la \VUgras{fortikaĵo}
\textsuperscript{(8)} kiu defendas la \VUgras{montpasejon}
\textsuperscript{(27)} apartigantan Francion de Hispanio tute memorigas tiujn
malnovajn kastelojn de la bordo de Rejno kiuj ŝajnas gvati \VUgras{rabaĵon} de
la supro de sia roko.

%%K t09-c1-05-1 p
5. --- Enorma \VUgras{aglo} \textsuperscript{(9)} kiu havas sian neston ne
malproksime de la \VUgras{forto} \textsuperscript{(8)} ofte altiras mian
atenton. Tiu \VUgras{rabobirdo}, kies \VUgras{beko} estas forta, ofte alportas
al siaj idoj ŝafidon aŭ kapridon kiujn ĝi tenas per siaj \VUgras{ungegoj}.
Tiel granda estas la grandeco de ĝiaj \VUgras{flugiloj} ke ĝi povas longtempe
glitflugi kun sia peza portaĵo.

%%K t09-tit-4 sub
\VUsaut{3.0mm}
\VUpk{10.2pt}{40}{La ekskursanto.}
\VUsaut{3.0mm}

%%K t09-c1-06-1 p
6. --- Tie la plej agrabla promeno estas la ekskurso al la kaskado
\textsuperscript{(10)} de « Kaŭ ». La \VUgras{akvofalo} kirle falas kaj poste
malaperas kiel bela \VUgras{rivereto} en \VUgras{arbaro de abioj}
\textsuperscript{(11)} situanta okcidente de la urbo. Ni supreniris tra ĉi tiu
arbaro ĝis la \VUgras{meteorologia observatorio} \textsuperscript{(12)} kaj de
la supro de la \VUgras{belvidejo}, ni ekvidis la tutan \VUgras{panoramon} de la
regiono. La \VUgras{pejzaĝo} estas mirinda, kaj la \VUgras{naturo} ŝajnas
malŝpari al la lando ĉiujn trezorojn kiujn ĝi disponas.

%%K t09-c1-07-1 p
7. --- Por malsupreniri, ni uzis la \VUgras{funikularon}
\textsuperscript{(13)} kaj ni haltis en malgranda \VUgras{stacidomo} apud la
\VUgras{ruinaĵoj} \textsuperscript{(14)} de malnova \VUgras{feŭda kastelo}
iam loĝita de la sinjoroj de Prats.

%%K t09-c1-08-1 p
8. --- Kompatinda kastelo! Kion kredeble ĝi pensas vidante malsupre la koketan
\VUgras{banurbon} \textsuperscript{(15)} kiu nun estas \VUgras{termala stacio}
laŭmoda?

%%K t09-c1-08-2 p
La \VUgras{klimato} estas tiel agrabla, la \VUgras{minerala akvo} tiel efika,
ke la \VUgras{kuracado} uzata en la \VUgras{sanatorio} \textsuperscript{(17)}
estas vera plezuro.

%%K t09-tit-5 sub
\VUsaut{3.0mm}
\VUpk{10.2pt}{40}{Agrabla termala urbo.}
\VUsaut{3.0mm}

%%K t09-c1-09-1 p
9. --- Cetere, oni faris ĉion por igi agrabla la restadon. Granda
\VUgras{viadukto} \textsuperscript{(16)} estis konstruita por ebligi fervojon;
la \VUgras{parko} \textsuperscript{(18)} de la \VUgras{termala establo}, en kiu
oni \VUgras{koncertas}, estas aranĝita en ĉarma maniero. Pluraj graciaj
« \VUgras{vilaoj} » \textsuperscript{(19)} estis konstruitaj ĉirkaŭ la kazino
\textsuperscript{(21)}, kaj \VUgras{ŝoseo} \textsuperscript{(22)} tre bone
tenata grande faciligas la interkomunikojn.

%%K t09-c1-10-1 p
10. --- Simpla \VUgras{taluso} \textsuperscript{(23)} apartigas la \VUgras{vojon}
\textsuperscript{(20)} de la \VUgras{valo} \textsuperscript{(25)} proprasence
dirita, kaj en kies fundo kuŝas gracia malgranda \VUgras{lago}
\textsuperscript{(26)} kiu nutras \VUgras{rivereton} \textsuperscript{(24)}
limpidan.

%%K t09-c1-11-1 p
11. --- Sur la alia flanko de la valo, \VUgras{ferminejo} \textsuperscript{(28)}
estas ekspluatata. Plurfoje mi iris vidi la \VUgras{ministojn}
malsuprenirantajn en la \VUgras{ŝakton}, kaj eĉ mi eniris la \VUgras{galerion}
en kiu ili pasigas sian tagon, eltirante la \VUgras{ercon} kiu enhavas la
\VUgras{metalon}.

%%K t09-c1-12-1 p
12. --- Grava \VUgras{fandejo} estis konstruita apud la minejo por tuj
utiligi la riĉaĵojn kiuj estas eltirataj el ĝi. La \VUgras{fornego}
\textsuperscript{(29)}, varmigita per la \VUgras{terkarbo} kiu abundas ĉi tie,
ebligas rapide kaj malkare akiri \VUgras{gisferon}.

%%K t09-c1-13-1 p
13. --- Ne malproksime de la minejo oni elfosas \VUgras{ŝtonminejon}
\textsuperscript{(30)}. Nek \VUgras{graniton}, nek \VUgras{porfiron}, nek eĉ
\VUgras{grejson} la maljuna \VUgras{ŝtonministo} detranĉas per sia
\VUgras{pioĉo}, sed simplan \VUgras{kvadron} por la konstruo de domoj.

%%K t09-c1-14-1 p
14. --- Unu el la plej belaj ornamaĵoj de tiu lando estas la \VUgras{abioj}
\textsuperscript{(31)}. Ĉie en la fundo de \VUgras{ravino}
\textsuperscript{(32)}, sur la bordo de \VUgras{torento} \textsuperscript{(33)},
de \VUgras{abismo} \textsuperscript{(34)} aŭ de krutaĵo, iliaj maldikaj pingloj
verde nuancas.

%%K t09-tit-6 sub
\VUsaut{3.0mm}
\VUpk{10.2pt}{40}{Piediro.}
\VUsaut{3.0mm}

%%K t09-c1-15-1 p
15. --- Mi profitas miajn feriojn por \VUgras{longe} promeni. La \VUgras{vojoj}
\textsuperscript{(35)} kiuj serpentumas sur la monto ebligas trakuri, sen
konsideri la distancon, plurajn \VUgras{kilometrojn} kaj ofte plurajn
\VUgras{kvaroblajn kilometrojn}.

%%K t09-c1-15-2 p
\VUgras{Limŝtonoj} \textsuperscript{(36)} kaj \VUgras{indikfostoj}
\textsuperscript{(37)} markas la vojojn sur kiuj ofte oni renkontas junan
\VUgras{montanon} \textsuperscript{(38)} kun \VUgras{dusako}
\textsuperscript{(40)} ĉe lia flanko, portante \VUgras{marmoton}
\textsuperscript{(39)}, aŭ kelkan \VUgras{ciganon} \textsuperscript{(41)} kies
\VUgras{feltoko} \textsuperscript{(42)} strange kontrastas kun la malnova
\VUgras{kapuĉo} \textsuperscript{(43)} ŝirita. Li kondukas per \VUgras{ĉeno
urson} \textsuperscript{(44)} dresitan.

%%K t09-tit-7 sub
\VUpk{10.2pt}{40}{Montsupreniro.}
\VUsaut{3.0mm}

%%K t09-c1-16-1 p
16. --- Preskaŭ ĉiutage mi iras kun \VUgras{gvidisto} \textsuperscript{(48)}
praktiki \VUgras{supreniron}, kaj mi estas tre fiera kiam, kun \VUgras{sako}
\textsuperscript{(47)} ĉe la dorso, kaj la « \VUgras{alpenstock} » en la mano,
kiel vera \VUgras{turisto} \textsuperscript{(46)} mi atakas la \VUgras{rokaĵojn}
\textsuperscript{(45)}.

%%K t09-c1-17-1 p
17. --- De la supro de monto, la loĝantoj de la ebenaĵo aspektas ne pli dikaj
ol formikoj; \VUgras{ŝafaro} \textsuperscript{(50)} paŝtanta en
\VUgras{paŝtejo} \textsuperscript{(49)} ŝajnas ludilo por infanoj.
\VUgras{Pino} \textsuperscript{(51)} aŭ ankaŭ \VUgras{ligna dometo}
\textsuperscript{(52)} aspektas apenaŭ kiel makulo sur la verdaĵo.

%%K t09-c1-18-1 p
18. --- Dum tiuj ekskursoj, ni ofte renkontas sur la \VUgras{vojeto}
\textsuperscript{(53)} \VUgras{ĉamĉasiston} \textsuperscript{(54)} ŝultre
portantan la beston kiun li ĵus mortigis.

%%K t09-c1-19-1 p
19. --- Mi lokis en mia herbario tre rimarkindan branĉon de \VUgras{filiko}
\textsuperscript{(55)} kaj belan rozkoloran branĉeton de \VUgras{eriko}
\textsuperscript{(57)} kiujn mi kolektis ĉe la enirejo de malgranda
\VUgras{groto} \textsuperscript{(56)} dum mia lasta promeno.

%%K t09-tit-8 sub
\VUsaut{3.0mm}
\VUcentre{12.0pt}{0}{\textit{Dua sceno.}}
\VUsaut{3.0mm}

%%K t09-tit-9 sub
\VUpk{11.4pt}{40}{La arbaro. --- La ĉasado.}
\VUsaut{3.0mm}

%%K t09-c2-01-1 p
1. --- Hieraŭ mi pasigis tagon delican en la arbaro. La diligenteco kiu regas
inter la bestetoj kaj la \VUgras{insektoj} \textsuperscript{(5)} de kiuj ĝi
estas loĝata tre interesis min.

%%K t09-c2-01-2 p
La \VUgras{araneo} \textsuperscript{(1)} kiu teksas sian \VUgras{reton}
\textsuperscript{(2)} inter du branĉoj por kapti la \VUgras{muŝon}
\textsuperscript{(3)} trairantan, la \VUgras{heliko} \textsuperscript{(4)}
penpene portanta sian konkon, la lukano esploranta pri sia kaptaĵo, la
diligenta formiko tre fervora apud la \VUgras{formikejo}
\textsuperscript{(6)}, tiuj ĉiuj malgranduloj iris, venis, laboris kun
eksterordinara diligenteco.

%%K t09-c2-02-1 p
2. --- \VUgras{Serpento} \textsuperscript{(7)} kiun mi unuavide opiniis
\VUgras{vipuro}, sed kiu estis nenio alia ol simpla \VUgras{kolubro}, estis
kaŝiĝinta sub arbotrunko, kaj tempe al tempo elirigis sian maldikan kapeton kiu
ĉiam ŝajnis preta por mordi. Antaŭ mi dika \VUgras{limako}
\textsuperscript{(8)} peze rampis post formanĝi unu el la bongustaj
\VUgras{fragetoj} \textsuperscript{(11)} kiuj abunde kreskas tie.

%%K t09-c2-03-1 p
3. --- La tero estis tapiŝita per \VUgras{arbaraj floroj}; \VUgras{primoloj}
\textsuperscript{(9)}, \VUgras{vinkoj} \textsuperscript{(10)}, kaj multaj
aliaj plantoj kiujn mi ne konis, floris inter \VUgras{herbotufoj}
\textsuperscript{(12)}.

%%K t09-c2-04-1 p
4. --- En tiu regiono, la \VUgras{trufo} estas preskaŭ tiel ordinara kiel la
\VUgras{fungo} \textsuperscript{(14)} kaj \VUgras{porkido}
\textsuperscript{(13)} apartenanta al arbaristo, manĝeme esploris ĉu ĝi
sukcesos trovi kelkajn el ili.

%%K t09-tit-10 sub
\VUsaut{3.0mm}
\VUpk{10.2pt}{40}{Ŝtelĉasado.}
\VUsaut{3.0mm}

%%K t09-c2-05-1 p
5. --- Mi ĉeestis la \VUgras{finon} de sceno kiu tre amuzis min. Per helpo de la
flarado de sia \VUgras{hundo baseto} \textsuperscript{(15)}, la
\VUgras{arbaristo} \textsuperscript{(16)} ĵus surprizis \VUgras{ŝtelĉasiston}
\textsuperscript{(22)}, en la momento kiam ĉi tiu preparis sin kunporti la
\VUgras{ĉasaĵon} \textsuperscript{(17)} kiun li estis mortiginta. Preferante
forlasi sian kaptaĵon ol lasi sin aresti, la ŝtelĉasisto estis forkurinta, kaj
\VUgras{leporo} \textsuperscript{(18)}, \VUgras{cervino} \textsuperscript{(19)},
\VUgras{kapreolo} \textsuperscript{(20)} kaj \VUgras{apro}
\textsuperscript{(21)} kuŝis sur la herbo ĝojigante la arbariston.

%%K t09-c2-06-1 p
6. --- La diverseco de la arboj kiujn enhavas la arbaro estas miriga. La
\VUgras{fagoj} \textsuperscript{(23)} kaj la \VUgras{betuloj}
\textsuperscript{(26)}, la \VUgras{pinoj}, kiuj produktas la \VUgras{rezinon}
\textsuperscript{(27)}, la \VUgras{kverkoj} \textsuperscript{(29)} kaj ankaŭ
multaj aliaj intermiksas sian branĉaron.

%%K t09-c2-06-2 p
La \VUgras{hedero} \textsuperscript{(24)}, kiu adheras al la trunko de la altaj
arboj, kolorigas la \VUgras{arbaregon} \textsuperscript{(25)} per verdaĵo
brilanta, kiu agrable unuigas sin al la nuanco pli klara kaj pale verda de la
\VUgras{musko} \textsuperscript{(31)} en kiu la \VUgras{verda pego}
\textsuperscript{(30)} serĉas insektojn.

%%K t09-tit-11 sub
\VUsaut{3.0mm}
\VUpk{10.2pt}{40}{Arbara pejzaĝo.}
\VUsaut{3.0mm}

%%K t09-c2-07-1 p
7. --- La malnovaj kverkoj ĉarmis min. Iliaj dikaj \VUgras{radikoj}
\textsuperscript{(32)} torditaj, duone elirintaj el la tero, donas al ili
belegan aspekton de forteco kaj vigoro, kaj mi demandas min kiel la
\VUgras{posedanto} \textsuperscript{(35)} povas rezigni pri faligi kaj liveri
ilin al la \VUgras{segilo} \textsuperscript{(34)} de la \VUgras{hakisto}
\textsuperscript{(33)}. La \VUgras{trunkoj} \textsuperscript{(36)}
kompatindaj, senigitaj je sia \VUgras{ŝelo} \textsuperscript{(37)} aŭ fenditaj
de la \VUgras{hakilo} \textsuperscript{(38)} de la \VUgras{taglaboristo}
\textsuperscript{(39)} ĉagrenigas min.

%%K t09-c2-08-1 p
8. --- Proksime, maljuna \VUgras{karbofaristo} \textsuperscript{(41)} preparis
per sia \VUgras{ŝovelilo amason} \textsuperscript{(42)} da karbo, kaj maljunulino
kunportis \VUgras{faskon} \textsuperscript{(40)} da \VUgras{ligno mortinta} el
branĉetoj de la faligitaj arboj.

%%K t09-tit-12 sub
\VUsaut{3.0mm}
\VUpk{10.2pt}{40}{Ĉasado.}
\VUsaut{3.0mm}

%%K t09-c2-09-1 p
9. --- Mi volis iri por ripozo dum kelkaj momentoj en la \VUgras{arbarista
domo} \textsuperscript{(46)} sed, kiam mi alvenis apud la \VUgras{lageto}
\textsuperscript{(43)}, sur kiu naĝas bela \VUgras{cigno}
\textsuperscript{(45)}, mi rimarkis, ke ĵusa \VUgras{inundo}
\textsuperscript{(44)} estis ŝanĝinta la vojon al vera \VUgras{marĉo}. Mi devis
deturni min, kaj, pasante apud la \VUgras{cedro} \textsuperscript{(49)}, la
\VUgras{poploj} \textsuperscript{(48)}, kaj la \VUgras{ulmo}
\textsuperscript{(47)} kiuj estas ĉe la bordo de la arbaro, mi vidis elveni el
\VUgras{arbetaĵo} \textsuperscript{(54)} belegan \VUgras{cervon}
\textsuperscript{(50)}, persekutatan de obstina \VUgras{(ĉas)hundaro}
\textsuperscript{(51)} kiun ekscitis ankaŭ la \VUgras{korno}
\textsuperscript{(53)} de la \VUgras{ĉasistoj rajdantaj} \textsuperscript{(52)}.
La kompatinda besto estis tute elĉerpita. Ĝi venis mortante fali je kelkaj
paŝoj apud mi.

%%K t09-c2-10-1 p
10. --- La filo de la arbaristo, kiun la bruado de la \VUgras{kurĉaso} estis
altirinta, eliris el la hejmo kaj ni duope revenis en la arbaron. Li estis
vidinta \VUgras{pigon} \textsuperscript{(56)} sur branĉo de malnova kverko. Li
opiniis, ke ĝia nesto kredeble estas kaŝita en la \VUgras{foliaro}
\textsuperscript{(58)} kaj li volis kapti la neston.

%%K t09-c2-10-2 p
Facilmova kiel \VUgras{sciuro} \textsuperscript{(61)}, li grimpis de
\VUgras{branĉo} \textsuperscript{(55)} al branĉo, sed li trovis nenion kaj
sukcesis nur ekflugigi malnovan \VUgras{kornikon} \textsuperscript{(60)}
sidantan sur \VUgras{branĉego} \textsuperscript{(57)}.
\VUsaut{4.0mm}
\VUfilet{22mm}
